\subsection{Riesz 空间的定义}

回顾一下，在集合 E 上，域 R 上的一个向量空间结构与一个序结构称为相容的，如果它们满足以下两条公理：

\((\mathrm{OVS}_{\mathrm{I}})\) 关系 \(x \leqslant y\) 蕴涵 \(x + z \leqslant y + z\) ，对一切 \(z \in \mathbb{E}\) 。

\((\mathrm{OVS}_{\mathrm{II}})\) 关系 \(x \geqslant 0\) 蕴涵 \(\lambda x \geqslant 0\) ，对每个标量 \(\lambda > 0\) 。

赋以这两个结构的空间 E 称为序向量空间 (TVS, II, §2, No. 5)。

公理 (OVS \(_{I}\) ) 表明 E 上的序结构与加群结构相容，换句话说，赋以这两个结构的 E 是一个序群 (A, VI, §1, No. 1)。

公理 (OVS \(_{\text{I}}\) ) 蕴涵关系 \(x \leqslant y\) 与 \(x + z \leqslant y + z\) 等价。类似地，由 (OVS \(_{\text{II}}\) ) 可得，对每个标量 \(\lambda > 0\) ，关系 \(x \leqslant y\) 与 \(\lambda x \leqslant \lambda y\) 等价，因为 \(\lambda^{-1} > 0\) ，从而关系 \(\lambda x \leqslant \lambda y\) 蕴涵 \(\lambda^{-1}(\lambda x) \leqslant \lambda^{-1}(\lambda y)\) 。因此可以说，在序向量空间中，平移以及比率 \(> 0\) 的位似都是序结构的自同构；这一事实也可以表述为：序在每次平移和每个比率 \(> 0\) 的位似下不变。此外，对称 \(x \mapsto -x\) 是 E 的序结构到相反序结构上的一个同构。

定义 1. --- 一个序向量空间称为 Riesz 空间（或格序向量空间）\(^{1}\) ，如果它的序结构是一个格序（也就是说，如果 E 的每对元素 x, y 都有上确界 \(\sup(x,y)\) 与下确界 \(\inf(x,y)\) ）。\(^{2}\)

例. 定义在任意集 A 上的一切实值函数所成的空间 \(\mathbf{R}^{\mathrm{A}}\) 是一个 Riesz 空间（对序关系 \(\ll x(t) \leqslant y(t)\) ，即对一切 \(t \in \mathrm{A} \gg\) ）；因为定义在 A 上的任何两个实值函数 \(x, y\) 都有上确界（相应地，下确界），它等于映射 \(t \mapsto \sup(x(t), y(t))\) （相应地，\(t \mapsto \inf(x(t), y(t))\) ）。

也可以这样说：Riesz 空间是赋以一个序结构的向量空间 E，使得一方面该结构与 E 的加群结构在 E 上定义一个格序群结构 (A, VI, §1, No. 9)，另一方面公理 (OVS \(_{II}\) ) 成立。

因此，格序群的所有性质都适用于 Riesz 空间；我们在此回顾其中的主要性质（参见 A, VI, §1, No. 9 至 No. 12），并指出由公理 (OVS \(_{II}\) ) 推出的那些结论。

首先回顾，我们记 \(x^{+} = \sup(x,0)\) （x 的正部分），\(x^{-} = (-x)^{+} = \sup(-x,0)\) （x 的负部分），\(|x| = \sup(x,-x)\) （x 的绝对值）；于是有 \(x = x^{+} - x^{-}\) 及 \(|x| = x^{+} + x^{-}\) ；在这里，这两个关系等价于

\[
x ^ {+} = \frac {1}{2} (| x | + x), \quad x ^ {-} = \frac {1}{2} (| x | - x).
\]

关系 \(x \leqslant y\) 等价于 \(\ll x^{+} \leqslant y^{+}\) 且 \(x^{-} \geqslant y^{-}\) 。对任意 \(x\) 与 \(y\) ，三角不等式成立：

\[
\vert x + y \vert \leqslant \vert x \vert + \vert y \vert .\tag{1}
\]

由序在每个比率 \textgreater{} 0 的位似下的不变性，

\[
\sup (\lambda x, \lambda y) = \lambda \sup (x, y) \quad \text { for   all } \lambda \geqslant 0.\tag{2}
\]

特别地，

\[
(\lambda x) ^ {+} = \lambda x ^ {+}, (\lambda x) ^ {-} = \lambda x ^ {-} \quad \text { for   all } \lambda \geqslant 0.\tag{3}
\]

另一方面，对 \(\lambda < 0\) ，我们有 \((\lambda x)^{+} = (-\lambda x)^{-} = |\lambda |x^{-}\) 及 \((\lambda x)^{-} = (-\lambda x)^{+} = |\lambda |x^{+}\) ；由此得出，对每个 \(\lambda \in \mathbf{R}\) 及每个 \(x\in \mathrm{E}\) ，

\[
\left| \lambda x \right| = \left| \lambda \right| \cdot \left| x \right|.\tag{4}
\]

序在平移下的不变性表明，对一切 \(z \in E\) ，

\[
\sup (x + z, y + z) = z + \sup (x, y),\tag{5}
\]

从而特别有，

\[
\sup (x, y) = x + (y - x) ^ {+} = \frac {1}{2} (x + y + | x - y |).\tag{6}
\]

我们有关系式

\[
\inf (x, y) = - \sup (- x, - y),\tag{7}
\]

\[
\sup (x, y) + \inf (x, y) = x + y.\tag{8}
\]

若 \(x, y, z\) 是 \(\geqslant 0\) 的，则 (A, VI, §1, No. 12, Prop. 11)

\[
\inf (x + y, z) \leqslant \inf (x, z) + \inf (y, z).\tag{9}
\]

若 A 与 B 是 E 的两个子集，且各有上确界，则 \(A + B\) 也有上确界，且

\[
\sup (A + B) = \sup A + \sup B.\tag{10}
\]

两个元素 \(x, y\) （属于 \(\mathrm{E}\) ）称为（相互）互奇异 \(^3\) 的，如果 \(\inf(|x|, |y|) = 0\) ；由 (8)，这一关系等价于 \(\sup(|x|, |y|) = |x| + |y|\) ，并且由 (6)，也等价于 \(||x| - |y|| = |x| + |y|\) ；0 是唯一与自身互奇异的元素；对每个 \(x \in \mathrm{E}\) ，\(x^+\) 与 \(x^-\) 互奇异，并且可以刻划为仅有的互奇异元素 \(y \geqslant 0\) 、\(z \geqslant 0\) （它们使 \(x = y - z\) ）。若 \(y\) 与 \(x\) 互奇异，则每个 \(z \in \mathrm{E}\) （它满足 \(|z| \leqslant |y|\) ）也与 \(x\) 互奇异。若 \(y\) 与 \(z\) 都与 \(x\) 互奇异，则由不等式 (9)，\(|y| + |z|\) 亦然；特别地，\(n|y|\) 与 \(x\) 互奇异对每个整数 \(n > 0\) 成立，由此推出 \(\lambda y\) 与 \(x\) 互奇异对每个标量 \(\lambda\) 成立，因为存在整数 \(n\) 使 \(|\lambda| \leqslant n\) ，从而 \(|\lambda y| \leqslant n|y|\) 。若一个子集 \(A\) （属于 \(E\) ）由与 \(x\) 互奇异的元素组成，且 \(A\) 有上确界，则该上确界也与 \(x\) 互奇异 (A, VI, §1, No. 12, Prop. 13 的推论)。

最后，我们有分解引理 (A, VI, §1, No. 10, Th. 1)：

若 \((x_{i})_{i\in \mathrm{I}}\), \((y_{j})_{j\in \mathrm{J}}\) 是 E 中元素 \(\geqslant 0\) 的两个有限序列，满足 \(\sum_{i\in \mathrm{I}}x_i = \sum_{j\in \mathrm{J}}y_j\) ，则存在一个有限序列 \((z_{ij})_{(i,j)\in \mathrm{I}\times \mathrm{J}}\) （由 E 中元素 \(\geqslant 0\) 组成），使得 \(x_{i} = \sum_{j\in \mathrm{J}}z_{ij}\) 对一切 \(i\in \mathrm{I}\) 成立，且 \(y_{j} = \sum_{i\in \mathrm{I}}z_{ij}\) 对一切 \(j\in \mathrm{J}\) 成立。

